% Part:sets-functions-relations
% Chapter: size-of-sets
% Section: enumerations

\documentclass[../../../include/open-logic-section]{subfiles}

\begin{document}

\olfileid[ar]{sfr}{siz}{enm}

\olsection{التعدادات ومجموعات \usetoken{S}{enumerable}}

\begin{editorial}
  يناقش هذا القسم تعدادات المجموعات، ويعرّفها بأنها دوال شاملة
  مجالها $\PosInt$. ويعرض الأمور بالتدرج مراعاةً للقراء ذوي الخلفية
  الرياضية المحدودة. وترد في \olref[enm-alt]{sec} نسخة بديلة أكثر إيجازًا،
  تعرّف التعدادات تعريفًا مختلفًا: بوصفها دوال تقابلية مع $\Nat$
  (أو مع مقطع ابتدائي).
\end{editorial}

\begin{explain}
سبق أن عرضنا أمثلة لمجموعات عن طريق سرد !!{element}sها. فلنناقش
بعبارات أعم كيف يمكننا سرد !!{element}s مجموعة ما ومتى يمكننا ذلك،
حتى إن كانت تلك المجموعة لا نهائية.
\end{explain}

\begin{defn}[التعداد، بصيغة غير رسمية]
بصيغة غير رسمية، \emph{تعداد} المجموعة~$A$ قائمة (قد تكون لا
نهائية) من !!{element}s~$A$، بحيث يظهر كل !!{element} من $A$
في القائمة عند موضع منتهٍ ما. وإذا كان للمجموعة $A$ تعداد، قيل إن
$A$ \emph{!!{enumerable}}.
\end{defn}

\begin{explain}
ثمة نقاط ينبغي ملاحظتها بشأن التعدادات:
\begin{enumerate}
\item لا نعدّ من التعدادات إلا القوائم التي لها بداية، والتي يكون
  لكل !!{element} فيها، عدا الأول، !!{element} واحد يسبقه مباشرة.
  وبعبارة أخرى، لا يفصل بين !!{element} القائمة الأول وأي
  !!{element} آخر إلا عدد منتهٍ من العناصر. ويعني هذا، على وجه
  الخصوص، أن لكل !!{element} في التعداد موضعًا منتهيًا: فموضع
  !!{element} الأول هو~$1$، وموضع الثاني~$2$، وهكذا.
\item قد تكون للمجموعة~$A$ تعدادات مختلفة تتباين بحسب ترتيب ظهور
  !!{element}sها: فالقائمة $4$، $1$، $25$، $16$،~$9$ تعدّد
  (مجموعة) المربعات الخمسة الأولى بالقدر نفسه من النجاح الذي تعدّدها
  به القائمة $1$، $4$، $9$، $16$،~$25$.
\item تظل التعدادات ذات التكرار تعدادات: فالقائمة $1$، $1$، $2$،
  $2$، $3$، $3$،~\dots{} تعدّد المجموعة نفسها التي تعدّدها القائمة
  $1$، $2$، $3$،~\dots{}.
\item إن الترتيب والتكرار \emph{مهمان بالفعل} حين نحدد تعدادًا: فيمكننا
  تعداد الأعداد الصحيحة الموجبة ابتداءً بـ$1$، $2$، $3$، $1$،
  \dots{}، لكن النمط يكون أوضح حين نعدّدها بالطريقة القياسية على
  الصورة $1$، $2$، $3$، $4$،~\dots
\item يجب أن تكون للتعدادات بداية: فالقائمة \dots، $3$، $2$، $1$
  ليست تعدادًا للأعداد الصحيحة الموجبة لأنها لا تضم !!{element} أولًا.
  ولكي ترى كيف ينتج هذا من التعريف غير الرسمي، فاسأل نفسك: «في أي موضع
  من القائمة يظهر العدد 76؟»
\item القائمة الآتية ليست تعدادًا للأعداد الصحيحة الموجبة:
  $1$، $3$، $5$، \dots، $2$، $4$، $6$، \dots\@ والمشكلة هي أن
  الأعداد الزوجية تقع في المواضع $\infty + 1$، $\infty + 2$، $\infty +
  3$، لا في مواضع منتهية.
\item المجموعة الخالية قابلة للتعداد: إذ تعدّدها القائمة الخالية!{}
\end{enumerate}
\end{explain}

\begin{prop}
  إذا كان للمجموعة $A$ تعداد، فلها تعداد بلا تكرار.
\end{prop}

\begin{proof}
  افترض أن للمجموعة $A$ تعدادًا $x_1$، $x_2$، \dots{} يكون فيه كل
  $x_i$ !!{element} من~$A$. ويمكننا إزالة التكرارات من تعداد
  بإزالة !!{element}s المكررة. فمثلًا، يمكننا تحويل التعداد إلى
  تعداد جديد نسرد فيه $x_i$، وهو !!a{element} من~$A$ لا يرد
  بين $x_1$، \dots، $x_{i-1}$، أو نحذف $x_i$ من القائمة إذا كان
  قد ورد من قبل بين $x_1$، \dots،~$x_{i-1}$.
\end{proof}

تبين الحجة الأخيرة أننا، لكي نحسن فهم التعدادات ومجموعات
!!{enumerable} ونبرهن نتائج عنها، نحتاج إلى تعريف أدق. ويقدمه التعريف
الآتي.

\begin{defn}[التعداد، بصيغة رسمية]
\emph{تعداد} مجموعة $A \neq \emptyset$ هو أي دالة
!!{surjective} $f \colon \PosInt \to A$.
\end{defn}

\begin{explain}
فلنتحقق من تكافؤ التعريف الرسمي والتعريف غير الرسمي الذي يستعمل قائمة
قد تكون لا نهائية. أولًا، تعدّد كل دالة !!{surjective} من $\PosInt$
إلى مجموعة~$A$ تلك المجموعة~$A$. فمثل هذه الدالة تحدد تعدادًا بالمعنى
غير الرسمي أعلاه، هو القائمة $f(1)$، $f(2)$، $f(3)$، \dots. وبما أن
$f$ !!{surjective}، فلا بد أن يكون كل !!{element} من~$A$ إحدى
قيم الدالة، أي~$f(n)$ لعدد ما~$n \in \PosInt$. ومن ثم يظهر كل !!{element} من
$A$ في موضع منتهٍ ما في القائمة. ولأن الدالة قد لا تكون
!!{injective}، فقد تتضمن القائمة تكرارًا، غير أن ذلك مقبول (كما
ذكرنا أعلاه).

ومن جهة أخرى، إذا أعطينا قائمة تعدّد جميع !!{element}s~$A$،
أمكننا تعريف دالة~!!a{surjective} $f\colon \PosInt \to A$ بأن نجعل
$f(n)$ هو !!{element} القائمة ذو الرتبة $n$، أو !!{element}ها
الأخير إذا لم يوجد فيها !!{element} ذو رتبة $n$. والحالة الوحيدة التي
لا ينتج فيها ذلك دالة~!!a{surjective} هي أن تكون $A$ خالية، ومن ثم
تكون القائمة خالية. وعلى هذا، تحدد كل قائمة غير خالية
دالة~!!a{surjective} $f\colon \PosInt \to A$.
\end{explain}

\begin{defn}
  \ollabel{defn:enumerable}
  تكون المجموعة~$A$ !!{enumerable} إذا وفقط إذا كانت خالية أو كان
  لها تعداد.
\end{defn}

\begin{ex}
الدالة التي تعدّد الأعداد الصحيحة الموجبة ($\PosInt$) هي ببساطة دالة
الهوية المعطاة بـ$f(n) = n$. والدالة التي تعدّد الأعداد الطبيعية
$\Nat$ هي الدالة $g(n) = n - 1$.
\end{ex}

\begin{ex}
الدالتان $f\colon \PosInt \to \PosInt$ و$g \colon \PosInt \to
\PosInt$ المعطاتان بـ
\begin{align*}
f(n) & = 2n \text{ و}\\
g(n) & = 2n - 1
\end{align*}
تعدّدان، على الترتيب، الأعداد الصحيحة الموجبة الزوجية والأعداد الصحيحة
الموجبة الفردية. غير أن أيًا منهما ليس تعدادًا لـ$\PosInt$، لأن كل
واحدة منهما ليست !!{surjective}.
\end{ex}

\begin{prob}
عرّف تعدادًا للمربعات الموجبة $1$، $4$، $9$، $16$، \dots
\end{prob}

\begin{ex}
تعدّد الدالة $f(n) = (-1)^{n} \lceil \frac{(n-1)}{2}\rceil$ (حيث
ترمز $\lceil x \rceil$ إلى \emph{دالة السقف}، التي تقرّب $x$ صعودًا
إلى أقرب عدد صحيح) مجموعة الأعداد الصحيحة~$\Int$. ولاحظ كيف
تولد $f$ قيم $\Int$ «بالقفز» جيئةً وذهابًا بين الأعداد الصحيحة الموجبة
والسالبة:
\[
\begin{array}{c c c c c c c c}
f(1) & f(2) & f(3) & f(4) & f(5) & f(6) & f(7) & \dots \\ \\
- \lceil \tfrac{0}{2} \rceil & \lceil \tfrac{1}{2}\rceil & - \lceil \tfrac{2}{2} \rceil & \lceil \tfrac{3}{2} \rceil & - \lceil \tfrac{4}{2} \rceil  & \lceil \tfrac{5}{2}
\rceil & - \lceil \tfrac{6}{2} \rceil & \dots \\ \\
0 & 1 & -1 & 2 & -2 & 3 & \dots
\end{array}
\]
ويمكنك أيضًا تصور $f$ على أنها معرفة بحسب الحالات كما يأتي:
\[
f(n) = \begin{cases}
  0 & \text{إذا كان $n = 1$}\\
  n/2 & \text{إذا كان $n$ زوجيًا}\\
  -(n-1)/2 & \text{إذا كان $n$ فرديًا و$>1$}
  \end{cases}
\]
\end{ex}

\begin{prob}
  أثبت أنه إذا كانت $A$ و$B$ !!{enumerable}، فإن $A \cup B$
  كذلك. ولإثبات ذلك، افترض وجود دالتين، كل منهما !!{surjective}، $f\colon \PosInt \to
  A$ و$g\colon \PosInt \to B$، وعرّف
  دالة~!!a{surjective}~$h\colon \PosInt \to A \cup B$ وبرهن أنها
  !!{surjective}. وانظر أيضًا في حالتي كون $A$ خالية أو~$B = \emptyset$.
\end{prob}

\begin{prob}
  أثبت أنه إذا كان $B \subseteq A$ وكانت $A$ !!{enumerable}،
  فإن~$B$ كذلك. ولإثبات ذلك، افترض وجود دالة~!!a{surjective} $f\colon \PosInt \to
  A$. وعرّف دالة~!!a{surjective}
  $g\colon \PosInt \to B$ وبرهن أنها !!{surjective}. ماذا يحدث إذا
  كان $B = \emptyset$؟
\end{prob}

\begin{prob}
  أثبت بالاستقراء على $n$ أنه إذا كانت $A_1$، $A_2$، \dots، $A_n$
  كلها !!{enumerable}، فإن $A_1 \cup \dots \cup A_n$ كذلك. ويجوز لك
  افتراض حقيقة أنه إذا كانت مجموعتان $A$ و~$B$ !!{enumerable}، فإن $A \cup
  B$ كذلك.
\end{prob}

مع أن من الأنسب، فيما يبدو، حين نسرد !!{element}s مجموعة، أن نبدأ
العد من !!{element} القائمة ذي الرتبة~$1$، فإن الرياضيين يحبذون استعمال
الأعداد الطبيعية~$\Nat$ لعد الأشياء. فهم يتحدثون عن !!{element}s
القائمة ذات الرتب $0$ و$1$ و$2$، وهكذا. وبالمثل، يمكننا تعريف
التعداد بأنه دالة~!!a{surjective} مجالها $\Nat$ ومجالها المقابل~$A$.
ومن الواضح أن التعريفين متكافئان.

\begin{prop}\ollabel{prop:enum-shift}
  توجد !!a{surjection} $f\colon \PosInt \to A$ إذا وفقط إذا وجدت
  !!a{surjection} $g\colon \Nat \to A$.
\end{prop}

\begin{proof}
  إذا أعطيت !!a{surjection} $f\colon \PosInt \to A$، أمكننا تعريف $g(n) =
  f(n+1)$ لكل $n \in \Nat$. ومن السهل التحقق من أن $g\colon \Nat
  \to A$ !!{surjective}. وبالعكس، إذا أعطيت !!a{surjection} $g\colon
  \Nat \to A$، فعرّف $f(n) = g(n-1)$.
\end{proof}

يعطينا هذا النتيجة الآتية:

\begin{cor}\ollabel{cor:enum-nat}
تكون المجموعة $A$ !!{enumerable} إذا وفقط إذا كانت خالية أو وجدت
دالة~!!a{surjective} $f\colon \Nat \to A$.
\end{cor}

ناقشنا أعلاه أن قائمة !!{element}s مجموعة~$A$ يمكن تحويلها إلى قائمة
بلا تكرار. ويصدق هذا أيضًا على التعدادات، لكن صياغة ذلك وبرهانه بدقة
أصعب قليلًا. فلا بد أن تكون كل دالة $f\colon \PosInt \to A$
معرفة عند كل $n \in
\PosInt$. وإذا لم يكن في~$A$ إلا عدد
منتهٍ من !!{element}s، فمن الواضح أنه لا يمكن أن تكون لدينا دالة معرفة
على عدد لا نهائي من !!{element}s~$\PosInt$، تتخذ جميع !!{element}s~$A$
قيمًا لها من غير أن تتخذ القيمة نفسها مرتين. وفي هذه الحالة، أي حين
تكون القائمة الخالية من التكرار منتهية، يجب أن نختار لـ$f$ مجالًا
آخر ليس فيه إلا عدد منتهٍ من !!{element}s. وعدم وجود تكرار يعني أن
$f$ يجب أن تكون !!{injective}. ولأنها !!{surjective} أيضًا، فنحن
نبحث عن !!a{bijection} بين مجموعة منتهية $\{1, \dots, n\}$ أو
$\PosInt$ وبين~$A$.

\begin{prop}\ollabel{prop:enum-bij}
إذا كانت $f\colon \PosInt \to A$ !!{surjective} (أي تعدادًا
لـ~$A$)، فثمة !!a{bijection} $g\colon Z \to A$، حيث $Z$
إما~$\PosInt$ وإما $\{1, \dots, n\}$ لعدد ما~$n \in \PosInt$.
\end{prop}

\begin{proof}
  نعرّف الدالة $g$ عوديًا: لنجعل $g(1) = f(1)$. وإذا كانت $g(i)$
  قد عُرّفت، فليكن $g(i+1)$ أول قيمة بين $f(1)$، $f(2)$، \dots{}
  لا ترد من قبل بين $g(1)$، \dots، $g(i)$، إن وجدت قيمة كهذه. فإذا
  كان لـ$A$ عدد قدره $n$ فقط من !!{element}s، كانت $g(1)$، \dots،
  $g(n)$ كلها معرفة، ومن ثم نكون قد عرّفنا دالة $g\colon \{1, \dots, n\}
  \to A$. وإذا كان لـ$A$ عدد لا نهائي من !!{element}s، فلا بد، لكل
  $i$، من وجود !!a{element} من~$A$ في التعداد $f(1)$، $f(2)$،
  \dots، لا يرد من قبل بين $g(1)$، \dots، $g(i)$. وفي هذه الحالة
  نكون قد عرّفنا دالة $g\colon \PosInt \to A$.

  الدالة $g$ !!{surjective}، لأن كل عنصر من~$A$ يرد بين $f(1)$،
  $f(2)$، \dots{} (بما أن $f$ !!{surjective})، ولذلك سيصبح في
  النهاية قيمةً لـ$g(i)$ عند عدد ما~$i$. وهي أيضًا !!{injective}،
  لأنه لو وجد $j < i$ بحيث $g(j) = g(i)$، لكانت $g(i)$ قد وردت
  من قبل بين $g(1)$، \dots، $g(i-1)$، خلافًا للطريقة التي عرّفنا
  بها~$g$.
\end{proof}

\begin{cor}\ollabel{cor:enum-nat-bij}
تكون المجموعة $A$ !!{enumerable} إذا وفقط إذا كانت خالية أو وجدت
!!a{bijection} $f\colon N \to A$، حيث إما $N = \Nat$ وإما
$N = \{0, \dots, n\}$ لعدد ما $n \in \Nat$.
\end{cor}

\begin{proof}
تكون $A$ !!{enumerable} إذا وفقط إذا كانت $A$ خالية أو وجدت دالة
!!a{surjective} $f\colon \PosInt \to A$. وبحسب
\olref{prop:enum-bij}، تصدق الحالة الأخيرة إذا وفقط إذا وجدت دالة
!!a{bijective}~$f\colon Z \to A$، حيث $Z =
\PosInt$ أو
$Z = \{1, \dots, n\}$ لعدد ما $n \in \PosInt$. وبالحجة نفسها
الواردة في برهان \olref{prop:enum-shift}، يصدق ذلك بدوره إذا وفقط إذا
وجدت !!a{bijection} $g\colon N \to A$، حيث إما $N = \Nat$ وإما
$N = \{0, \dots, n-1\}$.
\end{proof}

\begin{prob}
  بحسب \olref[sfr][siz][enm]{defn:enumerable}، تكون المجموعة $A$
  قابلة للتعداد إذا وفقط إذا كان $A = \emptyset$ أو وجدت دالة
  !!a{surjective} $f\colon
  \PosInt \to A$. ومن الممكن أيضًا أن
  نعرّف «مجموعة !!{enumerable}» تعريفًا دقيقًا كما يأتي: تكون مجموعة
  قابلة للتعداد إذا وفقط إذا وجدت دالة~!!a{injective}
  $g\colon A \to \PosInt$. أثبت تكافؤ التعريفين؛ أي أثبت أن دالة
  !!a{injective} $g\colon A \to \PosInt$ توجد إذا وفقط إذا كان
  $A = \emptyset$ أو وجدت دالة~!!a{surjective} $f\colon \PosInt \to A$.
\end{prob}

\end{document}
